\documentclass[11pt]{article}
\usepackage[margin=1in]{geometry}
\usepackage{amsmath,amssymb,amsthm,mathtools}
\usepackage{graphicx}
\usepackage{microtype}
\usepackage{xcolor}
\usepackage[hidelinks]{hyperref}

\newtheorem{theorem}{Theorem}
\newtheorem{lemma}{Lemma}
\newcommand{\C}{\mathbb C}
\newcommand{\Z}{\mathbb Z}
\newcommand{\B}{\mathcal B}
\newcommand{\K}{\mathcal K}
\newcommand{\Kz}{\mathcal K_0}
\newcommand{\Q}{\mathcal Q}
\newcommand{\dist}{\operatorname{dist}}

\title{A H\"older Counterexample to the Proposed
\texorpdfstring{$O(\log n/n)$}{O(log n/n)} Evolution--Trotter Rate}
\author{Autonomous proof attempt for arXiv:2002.04483}
\date{11 August 2026}

\begin{document}
\maketitle

\begin{center}
\fbox{\parbox{0.91\textwidth}{
\textbf{Status: counterexample likely valid, pending human review.}
For every $0<\beta<1$, a scalar commuting example satisfying assumptions
(A1)--(A5) of the source has error $\Omega(n^{-\beta})$ along dyadic $n$.
Hence the H\"older-time estimate (3.24) cannot in general be improved to
$O(\log n/n)$.  The separate Lipschitz question for (3.25) is not settled.
}}
\end{center}

\section{Source question}

The source is H.~Neidhardt, A.~Stephan, and V.~A.~Zagrebnov,
\emph{Operator-Norm Convergence of the Trotter Product Formula on Hilbert
and Banach Spaces: A Short Survey}, arXiv:2002.04483v1 (2020)~\cite{NSZsurvey}.
In Section~3.4, page~14, after Theorems~3.4--3.5, the authors ask whether
the rates (3.24) and (3.25) can be improved to $O(\log n/n)$.

\begin{figure}[h]
  \centering
  \includegraphics[width=0.96\textwidth]{figures/open_problem_crop.png}
  \caption{The exact open-question passage in arXiv:2002.04483v1, p.~14.}
\end{figure}

Theorem~3.4 assumes a H\"older exponent $\beta\in(\alpha,1)$ and gives
$O(n^{-(\beta-\alpha)})$.  The construction below disproves the proposed
$O(\log n/n)$ improvement for this H\"older branch.  It works in the scalar
model that the source itself identifies in Section~4.3 as a special case of
(A1)--(A5).

\section{Idea of the proof}

For a scalar time-dependent multiplication operator, evolution--Trotter error
is precisely the error of a uniform left- or right-endpoint Riemann sum.  We
build a nonnegative periodic H\"older function from dyadic tent waves.  On a
grid with $n=2^M$ points, every wave of frequency below $2^M$ has exactly the
correct discrete average, while every wave of frequency at least $2^M$
vanishes at all grid points.  The whole unresolved tail is therefore missed,
and its integral is exactly a positive constant times $n^{-\beta}$.

\section{Counterexample theorem}

Let $T=3$, $X=\C$, and let $D_0$ be the generator of the right-shift semigroup
on $L^2([0,T])$.  Fix $0<\alpha<\beta<1$, put $A=I$, and write
$\Kz=D_0+I$.  Define the one-periodic tent function and its
Takagi--Landsberg series by
\[
 \phi(x)=\dist(x,\Z),\qquad
 q(t)=\sum_{m=0}^{\infty}2^{-m\beta}\phi(2^m t).
\]
Let $B(t)=q(t)I$, let $\Q$ be multiplication by $q$, and let
$\K=D_0+I+\Q$.

\begin{theorem}\label{thm:main}
The family $(A,B(t))$ satisfies assumptions (A1)--(A5) of
\cite[Assumption~3.3]{NSZsurvey}, with the chosen $\alpha$ and $\beta$.
There is a constant $c_\beta>0$ such that, for every $M\ge1$ and
$n=2^M$,
\begin{align*}
 \sup_{\tau\ge0}\left\|
   (e^{-\tau\Q/n}e^{-\tau\Kz/n})^n-e^{-\tau\K}
 \right\| &\ge c_\beta n^{-\beta},\\
 \sup_{\tau\ge0}\left\|
   (e^{-\tau\Kz/n}e^{-\tau\Q/n})^n-e^{-\tau\K}
 \right\| &\ge c_\beta n^{-\beta}.
\end{align*}
In particular, neither error is $O(\log n/n)$.
\end{theorem}

\begin{proof}
We divide the proof into four short steps.

\emph{Step 1: $q$ is nonnegative, bounded, and $\beta$-H\"older.}
The first two assertions follow from $0\le\phi\le1/2$ and uniform
convergence.  Let $h=|t-s|\in(0,1]$ and choose $J\ge0$ so that
$2^Jh\le1<2^{J+1}h$.  Since $\phi$ is $1$-Lipschitz,
\begin{align*}
 |q(t)-q(s)|
 &\le h\sum_{m=0}^{J}2^{m(1-\beta)}
   +\sum_{m>J}2^{-m\beta}\\
 &\le C_\beta h^\beta.
\end{align*}
For $h>1$ the same estimate follows after enlarging $C_\beta$.  Thus
$q\in C^{0,\beta}([0,3])$.

\emph{Step 2: all source hypotheses hold.}
The scalar $A=I$ is invertible and generates a holomorphic contraction
semigroup.  Each $B(t)=q(t)I$ is bounded, nonnegative, and generates the
contraction $e^{-\tau q(t)}I$.  All domain, measurability, and adjoint
requirements in (A2)--(A4) are automatic.  Since every power of $A$ is $I$,
(A5) is exactly
\[
 \|A^{-1}(B(t)-B(s))A^{-\alpha}\|
 =|q(t)-q(s)|\le C_\beta|t-s|^\beta.
\]
The pair is also Trotter-stable because all factors are contractions.

\emph{Step 3: exact dyadic quadrature defect.}
Over any interval of length one,
\[
 Q_\beta:=\int_1^2q(t)\,dt
 =\frac14\sum_{m=0}^{\infty}2^{-m\beta}
 =\frac{1}{4(1-2^{-\beta})},
\]
because $\int_0^1\phi(u)\,du=1/4$.  Take $n=2^M$.  For $m<M$, the
grid average of $\phi(2^m t)$ on $[1,2]$ is $1/4$: after removing repeated
periods it is
\[
 \frac1L\sum_{j=0}^{L-1}\phi(j/L)=\frac14,
 \qquad L=2^{M-m}\ge2.
\]
For $m\ge M$, every grid value is zero.  Hence both endpoint sums agree and
\begin{align*}
 R_M
 &:=\frac1n\sum_{k=0}^{n-1}q(1+k/n)
  =\frac1n\sum_{k=1}^{n}q(1+k/n)\\
 &=\frac14\sum_{m=0}^{M-1}2^{-m\beta},
\end{align*}
so that
\begin{equation}\label{eq:defect}
 \delta_M:=Q_\beta-R_M
 =\frac{n^{-\beta}}{4(1-2^{-\beta})}.
\end{equation}

\emph{Step 4: transfer the defect to operator norm.}
The exact propagator for $D_0+I+\Q$ on the interval $[s,t]$ is the scalar
multiplier
\[
 U(t,s)=\exp\left(-(t-s)-\int_s^tq(y)\,dy\right).
\]
The two product orders replace the integral by the left- and right-endpoint
sums, respectively.  At the interior point $(t,s)=(2,1)$ those sums both
equal $R_M$.  Proposition~3.8 of the source identifies the evolution-semigroup
operator norm with the essential supremum of the propagator error.  The
multipliers are continuous, so this essential supremum is at least their
value at $(2,1)$.  By \eqref{eq:defect}, for either product order it is at
least
\begin{align*}
 e^{-1-R_M}-e^{-1-Q_\beta}
 &=e^{-1-Q_\beta}(e^{\delta_M}-1)\\
 &\ge \frac{e^{-1-Q_\beta}}{4(1-2^{-\beta})}\,n^{-\beta}.
\end{align*}
This proves the asserted lower bound.  Finally,
$n^{-\beta}/(\log n/n)=n^{1-\beta}/\log n\to\infty$, so an
$O(\log n/n)$ estimate is impossible.
\end{proof}

\section{Verification and scope}

The script \texttt{code/check\_dyadic\_identity.py} checks finite truncations
of the exact quadrature identity, the constancy of $n^\beta\delta_M$, and the
semigroup lower bound.  It is verification only; the proof above is exact.

The result gives a complete negative answer to the (A5), $0<\beta<1$ half of
the source's proposed $O(\log n/n)$ improvement.  It does not address the
separate Lipschitz assumption (A6) or whether the general upper rate
$O(n^{-(\beta-\alpha)})$ can be sharpened to the scalar-optimal
$O(n^{-\beta})$.

\section{Bounded novelty check}

We searched the run indexes and exact arXiv id; the full sources of
arXiv:1612.06147, 1703.09536, 1901.02205, and 2002.04483; and arXiv API
queries through 11 August 2026 for combinations of ``Trotter product'',
``H\"older'', ``non-autonomous/evolution'', ``sharp'', and ``Riemann sums''.
The cited scalar paper proves the upper $O(n^{-\beta})$ estimate and arbitrary
slowness for merely continuous functions, but no H\"older lower example or
application disproving the survey's open $O(\log n/n)$ question was found.
The novelty claim is necessarily bounded; the dyadic construction is an
elementary sharp quadrature example and may be known outside the searched
Trotter literature.

\paragraph{Human-review recommendation.}
Check the source's quantifier in the plural open question and the transfer from
the interior propagator value to the essential supremum.  If confirmed, accept
this as a counterexample to the H\"older branch, explicitly retaining the A6
limitation.

\begin{thebibliography}{9}
\bibitem{NSZsurvey}
H.~Neidhardt, A.~Stephan, and V.~A.~Zagrebnov,
\emph{Operator-Norm Convergence of the Trotter Product Formula on Hilbert and
Banach Spaces: A Short Survey}, arXiv:2002.04483v1, 2020.

\bibitem{NSZremarks}
H.~Neidhardt, A.~Stephan, and V.~A.~Zagrebnov,
\emph{Remarks on the operator-norm convergence of the Trotter product
formula}, arXiv:1703.09536v1, 2017.
\end{thebibliography}

\end{document}

